\begin{proposition}[Canonical rings and rational topology of the covers]\hypertarget{tag-01S9}{}\label{inv:B3.3}
\StatementRecord{inv:B3.3}{proved}{01S9}{\SeriesRef{inv:B3.1}, \SeriesRef{inv:B2.2}, \SeriesRef{inv:B3.2}}

For $X$ and $Y$ as above,
\[
K_Y=\pi^*K_X=\OO_Y(-2S'_2),\quad
\bigoplus_{m\ge0}H^0(Y,-mK_Y)=\C[U,V],\quad
\deg(U,V)=(1,2).
\]
The same ring holds for $X$.
Over $\C$, the Betti numbers of $Y$ are
$(1,0,n-1,0,n-1,0,1)$; its Fr\"olicher sequence degenerates at $E_2$.
For every nontrivial character of the deck group, the character
part of Dolbeault cohomology consists of one class in $(1,1)$ and
one in $(2,2)$ and no others.
\end{proposition}


\begin{proof}
\begin{enumerate}[label=\arabic*.,leftmargin=*]
\item The canonical formula pulls back through an \'etale map, and
$4S'_2=g^*p_2$. \ref{inv:B3.1} gives $g_*\OO_Y=\OO_B$.
The valuation argument in \ref{inv:B2.2} therefore applies unchanged and proves
the ring statement with all multiplication maps.

\item The comparisons of \ref{inv:B3.1} are pullback maps and identify the logarithmic
cohomology bundles with the invariant part; the deck group acts
trivially on those bundles. At the cusp its action on the $n$
normalization components is the regular permutation representation.
The quotient by the diagonal has each nontrivial character once.
Thus the first-row additional skyscraper in \ref{inv:B3.2} has precisely these
characters, and duality gives the conjugate characters in its
second-row additional skyscraper. All other cohomology is invariant
and equal to that of $X$. Coherent Leray now proves the assertion
about the Dolbeault character spaces.

\item The de Rham complex and its Fr\"olicher sequence split under the
idempotents of the finite deck group. A nontrivial character has just
the two stated Dolbeault entries. No differential can connect them
or connect either to another nonzero entry, so both survive.
The invariant summand is the de Rham complex of $X$ by finite
\'etale descent and averaging. Its rational topology is computed by the following attachment maps. In the remainder of this proof
all homology and cohomology coefficients are $\Q$.

\item Let $N$ be a small closed cusp tube, $M=\partial N$, and
$U=X\setminus\operatorname{int}N$. The space $U$ is covered by
two pieces retracting onto $S_1,S_2$, whose intersection retracts
onto a regular fibre $F$. The retractions are obtained from the
equivariant smooth torus trivializations used in \ref{inv:B3.1}. The maps
$F\to S_j$ are the central finite covering maps; translation of
their origins induces the identity on their rational cohomology.
Finite-cover transfer identifies their pullbacks with the invariant
subspaces for $A_j$. Put
\[
 \psi_1=2u+w+3\delta,\qquad \psi_2=u+w+2\delta,
 \qquad \eta=u\wedge w+6\gamma\wedge\delta.
\]
Direct substitution in $A_1,A_2$ shows that the invariant
degree-one spaces are $\langle\gamma,\psi_j\rangle$ and the
invariant degree-two spaces are
$\langle\gamma\wedge\psi_j,\eta\rangle$.
Their dimensions are two: the nontrivial eigenvalues of $A_j$ on
the complementary plane are the inverse primitive third or fourth
roots, so the exterior-square fixed space has exactly the two
displayed generators.
In the transfer-normalized dual homology bases, the two
Mayer--Vietoris maps are therefore the following matrices:
\[
 \alpha_1=\begin{pmatrix}
 1&0&0&0\\0&2&1&3\\-1&0&0&0\\0&-1&-1&-2
 \end{pmatrix},\qquad
 \alpha_2=\begin{pmatrix}
 2&1&3&0&0&0\\0&0&6&1&0&0\\
 -1&-1&-2&0&0&0\\0&0&-6&-1&0&0
 \end{pmatrix}.
\]
The second matrix uses the ordered wedge basis
$(\hat\gamma\hat u,\hat\gamma\hat w,\hat\gamma\hat\delta,
\hat u\hat w,\hat u\hat\delta,\hat w\hat\delta)$.
Both matrices have rank three, and
$\ker\alpha_1=\Q\nu_0$, with
$\nu_0=-\hat u-\hat w+\hat\delta$.
Consequently
\begin{equation}
 H_1(U)=\Q,\qquad
 0\longrightarrow\operatorname{coker}\alpha_2=\Q
 \longrightarrow H_2(U)
 \xrightarrow{\partial}\Q\nu_0\longrightarrow0.
\label{pub:I:eq:B-1}
\end{equation}
The generator $\Theta\in H^1(U)$ is normalized by
$\Theta|_F=\gamma$. The finite affine generators have
$3\bar g_1=l_1c$, $4\bar g_2=l_2c$ in rational homology, where
$c$ is the image of $\hat\gamma$ and $\Theta(c)=1$.
For the toric meridian $a_0=t_\mu y^{-1}x^{-1}$ this gives the
explicit collar value
\begin{equation}
 \Theta(j_*a_0)=l_0-l_1/3-l_2/4=p/12.
\label{pub:I:eq:B-2}
\end{equation}

\item Lemma~\ref{I:lem:cusp-tube-contraction} gives the strong
deformation retraction $N\to W$, fixing $W$ pointwise.

\item The constant-sheaf normalization resolution of $W$ has strata
$\mathrm{dP}_6$, three copies of $\Pone$, and two points.
Its row of degree-zero cohomology is
\[
 \Q\xrightarrow{0}\Q^3
 \xrightarrow{(a,b,c)\mapsto(a-b+c,a-b+c)}\Q^2.
\]
The strata have no degree-one cohomology. This resolution therefore
gives $H_1(N)=H_1(W)=\Q^2$, $H_4(N)=\Q[W]$ and $H_5(N)=0$;
it also gives $e(W)=6-3\cdot2+2=2$.
Poincar\'e--Lefschetz duality for the compact smooth manifold with
boundary $(N,M)$ now shows that $H_1(M)\to H_1(N)$ and
$H_2(M)\to H_2(N)$ are surjective. For degree one the next
relative group is $H_1(N,M)=H^5(N)=0$.
For degree two the map to $H_2(N,M)=H^4(N)$ is the intersection
pairing with $[W]$, which is zero because $W$ is the divisor of
the global function $t_c$ on $N$. The Wang sequence for the
boundary mapping torus, with $\rank(M_0-I)=2$, gives
$\dim H_1(M)=3$. Its kernel to $H_1(N)$ is exactly
$\Q a_0$: this meridian is nonzero under projection to the base
circle and bounds its normal disc in $N$.

\item We exhibit two classes in
$K_2=\ker(H_2(M)\to H_2(N))$.
For $\nu=\hat w,\hat\delta$, let $\xi_\nu$ be the two-torus
obtained by translating the toric meridian around the angular
circle $S^1_\nu$. In the chart $t_c=z_0z_1z_2$ corresponding to
the rays $(0,1),(e_1,1),(e_2,1)$, these are the boundaries of
$D^2_{z_0}\times S^1_{z_i}$ with the other coordinate fixed
and the radii chosen so that $|t_c|=\epsilon$ on the boundary.
They therefore bound in $N$ with the specified toric meridian.
Cutting the meridian into the two arcs of the cover defining
\eqref{pub:I:eq:B-1}, the two cylinders have end circles differing by
$(I-A_1^{-1})\nu$. With this orientation of the Mayer--Vietoris
boundary, matrix multiplication gives
\begin{equation}
 \partial j_*\xi_{\hat w}=-\nu_0,\qquad
 \partial j_*\xi_{\hat\delta}=0.
\label{pub:I:eq:B-3}
\end{equation}
The affine translation of either endpoint has no effect on its
homology class, which explains why this boundary formula is the
same for every translation twist.

\item Choose $\mathcal A\in H^2(U)$ restricting to $\eta$ on $F$;
it exists by Mayer--Vietoris because the two transfer-normalized
classes $\eta$ have equal restrictions. Fibrewise translation by
$S^1_{\hat\delta}$ defines a global smooth action on $U$:
$A_j\hat\delta=\hat\delta$, and all affine translations commute
with this circle. Slant product of its action map with the circle
defines $\iota:H^2(U)\to H^1(U)$. Naturality on $F$ gives
$\iota\mathcal A|_F=\iota_{\hat\delta}\eta=-6\gamma$.
The restriction $H^1(U)\to H^1(F)$ is injective by \eqref{pub:I:eq:B-1}, so
$\iota\mathcal A=-6\Theta$. By the definition of slant product,
\begin{equation}
 \langle\mathcal A,j_*\xi_{\hat\delta}\rangle
 =\langle\iota\mathcal A,j_*a_0\rangle=-p/2.
\label{pub:I:eq:B-4}
\end{equation}
We may add the class pulled back from $(\Q\nu_0)^*$ under
the boundary map in \eqref{pub:I:eq:B-1} so that
$\langle\mathcal A,j_*\xi_{\hat w}\rangle=0$;
this does not change its restriction to $F$ or \eqref{pub:I:eq:B-4}.
In the coordinates on $H_2(U)$ given by $\partial/\nu_0$ and
this normalized pairing, the two attachment cycles have
matrix
\begin{equation}
 \bigl[j_*\xi_{\hat w}\ \ j_*\xi_{\hat\delta}\bigr]
 =\begin{pmatrix}-1&0\\0&-p/2\end{pmatrix}.
\label{pub:I:eq:B-5}
\end{equation}
It is invertible over $\Q$ because admissibility implies
$p\ne0$. Thus the image of $K_2$ already fills $H_2(U)$.
Likewise \eqref{pub:I:eq:B-2} says the image of $\ker(H_1(M)\to H_1(N))$
is nonzero in $H_1(U)$. Using the surjectivity to $N$ just proved,
Mayer--Vietoris for $X=U\cup N$ now gives $H_1(X)=H_2(X)=0$.
Every other fibre has Euler number zero, so additivity over the
three special fibres and the smooth torus bundle gives $e(X)=2$.
Poincar\'e duality then gives $b_4=b_5=0$ and $b_3=0$.
This proves the asserted rational homology of $X$ for all admissible
twists; the only changing determinant in \eqref{pub:I:eq:B-2}--\eqref{pub:I:eq:B-5} is a nonzero
constant multiple of $p$. 

\item Proposition~\ref{inv:C1.1}, applied to this rational topology and
the same Hodge table, gives degeneration of the invariant summand at $E_2$.
Summing it with the nontrivial character spaces proves the Betti
numbers and the degeneration statement for $Y$. In particular
$e(Y)=2n$, as also follows from the covering degree.
\end{enumerate}
\end{proof}